\documentclass{article}
\usepackage[margin=2cm]{geometry}
\usepackage{amsmath}
\usepackage{graphicx}
\usepackage{booktabs}
\usepackage[utf8]{inputenc}
\begin{document}
allows us to conclude that this is always the case for real initial conditions \(\phi_{0}^{0}\). Thus we have a one parameter family of real smooth solutions, labeled by the IR parameter \(\phi_{0}^{0}\). With this in mind, we may choose any value of \(\phi_{0}^{0}\) and solve the BPS equations in numerically. In Figure , we plot the solutions obtained for the following choices of initial condition: \(\phi_{0}^{0}=\{0.25,0.5,1,1.5,2\}\). In order to get smooth solutions for \(u>0\), we must take \(\eta=-1\). It is straightforward to verify that the resulting solutions are completely smooth and have the expected vanishing of \(e^{2 f}\) at the origin, implying that the spacetime smoothly pinches off. Furthermore, \(e^{2 f} / e^{2 u}\) is seen to asymptote to a constant, which we denote by \(e^{2 f_{k}}\).
\subsection*{0.1 UV asymptotic expansions}
As in the holographic Janus solutions in Lorentzian signature, the BPS equations may also be used to obtain the UV asymptotic behavior of the solutions. To do so, we begin by defining an asymptotic coordinate \(z=e^{-u}\), where the asymptotic \(S^{5}\) boundary is reached by taking \(u \rightarrow \infty\). Consequently, an asymptotic expansion is an expansion around \(z=0\). The coefficients in the UV expansions of the non-zero fields may now be solved for order-by-order using the BPS equations. One finds explicitly that all coefficients are determined in terms of only three independent parameters \(\alpha, \beta\), and \(f_{k}\), in accord with the fact that there are three independent first-order differential equations. The first few terms in the expansions are
\[
\begin{aligned}
f(z)= & -\log z+f_{k}-\left(\frac{1}{4} e^{-2 f_{k}}+\frac{1}{16} \alpha^{2}\right) z^{2}+O\left(z^{4}\right) \\
& \sigma(z)=\frac{3}{8} \alpha^{2} z^{2}+\frac{1}{4} e^{f_{k}} \alpha \beta z^{3}+O\left(z^{4}\right) \\
\phi^{0}(z)= & \alpha z-\left(\frac{5}{4} \alpha e^{-2 f_{k}}+\frac{23}{48} \alpha^{3}\right) z^{3}+O\left(z^{4}\right) \\
& \phi^{3}(z)=e^{-f_{k}} \alpha z^{2}+\beta z^{3}+O\left(z^{4}\right)
\end{aligned}
\]
We have obtained the expansions up to \(O\left(z^{8}\right)\), but we display only the first few terms here.
\section*{1 Holographic sphere free energy}
The goal of this section is to obtain the holographic free energy, i.e. the renormalized on-shell action. We begin by writing the full action,
\[
\begin{array}{c}
S=S_{6 \mathrm{D}}+S_{\mathrm{GH}} \\
S_{6 \mathrm{D}}=\int d u d^{5} x \sqrt{G} \mathcal{L} \\
S_{\mathrm{GH}}=-\frac{1}{2} \int d^{5} x \sqrt{\gamma} \mathcal{K}
\end{array}
\]
where \(S_{6 \mathrm{D}}\) is the six-dimensional Euclidean action given in and \(S_{\mathrm{GH}}\) is the Gibbons-Hawking term. The \(\gamma\) appearing in \(S_{\mathrm{GH}}\) is the determinant of the induced metric on the boundary
\end{document}
